% !TeX program = pdflatex
% =============================================================================
% Aufgabenblatt 4: Wechselwirkung
% UZH Praktikum I -- Sara Romer Urban, Kantonsschule im Lee, HS2026
% Klasse: 2e (16.09.2026)
%
% Reine Uebungsaufgaben zur Wechselwirkung, kein einleitender Erklaerkasten
% (siehe institutionen/UZH-Praktikum-I/CLAUDE.md). Folgt lernziele.md,
% Abschnitt "Aufteilung Arbeitsblatt/Aufgabenblatt": (1) Zeichenaufgabe 2 --
% Objekt auf einer Oberflaeche, strukturell identisch zur Federwaage-Aufgabe
% im Arbeitsblatt, jetzt selbstaendig und als eigene TikZ-Skizze statt
% wiederverwendeter Normalkraft-Bilder (LZ2, zweiter Teil) -- (2) die
% restlichen zwei Rueckstossprinzip-Beispiele: Flugzeug, Person die einen
% Stein aus einem Boot wirft (LZ6) -- (3) Log-Rolling als Reflexionsfrage,
% weshalb das Kraeftepaar hier weniger klar isolierbar ist (LZ6) -- (4)
% optionaler Hinweis auf PhET "Forces and Motion: Basics" (Tauziehen-Screen,
% LZ4) zur Selbstkontrolle -- (5) optionale Zusatzaufgabe: Federwaage-
% Massen-Experiment mit anderen Werten (LZ3, produktives Ueben).
%
% Sprinter und Ruderboot (urspruenglich Teil von (2)) wurden gestrichen:
% beide stellten denselben kognitiven Schritt wie Flugzeug/Stein-aus-Boot
% ohne zusaetzliche Tiefe dar (reine "Wechselwirkungspaar benennen"-
% Wiederholung), Sprinter zudem strikt flacher als die bereits im
% Arbeitsblatt behandelte Skateboard-Variante 3 (Wand), die dieselbe
% Push-off-Situation schon mit Begruendung ueber die effektive Masse
% vertieft (siehe redundanz-tiefe-reviewer).
%
% Quellen: lernziele.md (dieser Ordner).
% =============================================================================
\documentclass[addpoints,11pt,a4paper]{exam}

\usepackage[T1]{fontenc}
\usepackage[utf8]{inputenc}
\usepackage[ngerman]{babel}
\usepackage{lmodern}
\usepackage[a4paper,margin=2.2cm,headheight=14pt]{geometry}
\usepackage{amsmath}
\usepackage{siunitx}
% Hausstil: zusammengesetzte Einheiten immer als echter Bruch (\frac), nicht
% als Inline-Schraegstrich oder \tfrac -- gilt global fuer jedes \si/\SI.
\sisetup{per-mode=fraction,fraction-command=\frac}
\usepackage{array,booktabs,tabularx,multirow}
\usepackage{enumitem}
\usepackage{xcolor}
\usepackage{colortbl}
\usepackage{tikz}
\usepackage{physics}
\usepackage[most]{tcolorbox}
\usepackage{lastpage}
\usepackage{graphicx}
\usepackage{hyperref}

\usetikzlibrary{arrows.meta,calc,angles,quotes,decorations.pathmorphing}

\usepackage{setspace}
\usepackage{needspace}
\setlength{\parindent}{0pt}
\setlength{\parskip}{0.6em}
\setstretch{1.08}
\setlist[itemize]{itemsep=0.4em}

% -----------------------------
% Farben (Corporate-Farbschema Physik KFR)
% -----------------------------
\definecolor{titleblue}{HTML}{183A6B}
\definecolor{accentblue}{HTML}{2E6F9E}
\definecolor{lightblue}{HTML}{EAF4FB}
\definecolor{lightgray}{HTML}{F4F6F8}
\definecolor{darkgray}{HTML}{4A4A4A}
\definecolor{accentorange}{HTML}{D9720C}
\definecolor{accentpurple}{HTML}{7B2D8E}
% Hinweisbox: Orange (accentorange), fuer wichtige Klarstellungen.
\definecolor{lightorange}{HTML}{FCEEE0}
% Formelbox: Violett (accentpurple), fuer die zentralen Merksaetze.
\definecolor{lightpurple}{HTML}{F3EAF7}

% Links (hyperref): farblich ans Corporate-Farbschema angeglichen.
\hypersetup{
  colorlinks=true,
  urlcolor=accentblue,
  linkcolor=accentblue,
  pdftitle={Aufgabenblatt 4: Wechselwirkung},
  pdfauthor={Loris Keller}
}

% -----------------------------
% Spaltentypen
% -----------------------------
\newcolumntype{Y}{>{\centering\arraybackslash}X}
\newcolumntype{P}[1]{>{\raggedright\arraybackslash}p{#1}}

% -----------------------------
% Boxen
% -----------------------------
\tcbset{before skip=10pt, after skip=10pt}
\newtcolorbox{titlebox}{
  enhanced,
  colback=titleblue,
  colframe=titleblue,
  boxrule=0pt,
  arc=3mm,
  left=3mm,right=3mm,top=3mm,bottom=3mm
}

\newlength{\aufgabenindent}
\settowidth{\aufgabenindent}{10.\hskip\labelsep}

% Praktikum I: Lernziele/Einleitung/Theorie/Aufgaben werden NICHT mehr
% farbig gerahmt (nur Formel-, Hinweis- und Antwortbox bleiben Boxen) --
% siehe institutionen/UZH-Praktikum-I/CLAUDE.md, Abschnitt "Boxen-Layout".
\newenvironment{infobox}{%
  \par\addvspace{10pt}\noindent
}{%
  \par\addvspace{10pt}
}

\newenvironment{theoriebox}[1][Theorie]{%
  \par\addvspace{10pt}\noindent\textbf{#1}\par\smallskip\noindent
}{%
  \par\addvspace{10pt}
}

% taskbox/groupbox stehen in questions VOR dem ersten \question (Bilder,
% Fliesstext) -- exam.cls' questions-Liste braucht dafuer eine echte,
% eigenstaendige Box (sonst "Something's wrong--perhaps a missing \item",
% weil z.B. ein \begin{center} vor dem ersten \item der Aussenliste steht).
% Deshalb bleiben es tcolorboxen, nur unsichtbar gemacht (keine Farbe/Rahmen,
% kein Padding) statt sie durch reine Absaetze zu ersetzen.
\newtcolorbox{taskbox}[1][]{
  enhanced, breakable,
  colback=white, colframe=white, boxrule=0pt,
  left=0mm,right=0mm,top=0mm,bottom=0mm,
  before skip=10pt, after skip=10pt,
  before upper={%
    \def\tmpboxtitle{#1}%
    \ifx\tmpboxtitle\empty\else\noindent\textbf{#1}\par\smallskip\fi
  }
}

\newtcolorbox{groupbox}[1][Gruppenarbeit]{
  enhanced, breakable,
  colback=white, colframe=white, boxrule=0pt,
  left=0mm,right=0mm,top=0mm,bottom=0mm,
  before skip=10pt, after skip=10pt,
  before upper={\noindent\textbf{#1}\par\smallskip}
}

% beispielbox: fuer die Einleitung -- wie die anderen Inhaltsboxen ungerahmt.
\newenvironment{beispielbox}[1][Beispiel]{%
  \par\addvspace{10pt}\noindent\textbf{#1}\par\smallskip\noindent
}{%
  \par\addvspace{10pt}
}

% hinweisbox: Orange, fuer wichtige Klarstellungen/Verwechslungswarnungen.
\newtcolorbox{hinweisbox}[1][Hinweis]{
  enhanced,
  breakable,
  colback=lightorange,
  colframe=accentorange,
  title={#1},
  fonttitle=\bfseries,
  boxrule=0.7pt,
  arc=2mm,
  left=5mm,right=5mm,top=2mm,bottom=2mm
}

% formelbox: Violett, um die zentralen Merksaetze dieser Einheit hervorzuheben.
\newtcolorbox{formelbox}[1][Merksatz]{
  enhanced,
  breakable,
  colback=lightpurple,
  colframe=accentpurple,
  title={#1},
  fonttitle=\bfseries,
  boxrule=0.9pt,
  arc=2mm,
  left=5mm,right=5mm,top=2mm,bottom=2mm
}

\newtcolorbox{answerbox}{
  breakable,
  colback=white,
  colframe=gray!55,
  boxrule=0.5pt,
  arc=1.5mm,
  left=2mm,right=2mm,top=2mm,bottom=2mm
}

% Marke: kleines farbiges Inline-Label fuer Teilschritte/Ergebnis-Callouts.
\newtcbox{\Marke}[1][accentpurple]{
  nobeforeafter,
  colback=#1,
  colframe=#1,
  coltext=white,
  boxrule=0pt,
  arc=2pt,
  left=5pt,right=5pt,top=2pt,bottom=2pt,
  fontupper=\bfseries\sffamily\small
}

% Bild-Helfer: zentriertes Bild mit spuerbarem Abstand zum umgebenden Text
% (statt nur des normalen Absatzabstands). #1 = Breite, #2 = Dateipfad.
\newcommand{\Bild}[2]{%
  \par\vspace{0.6cm}
  \begin{center}
    \includegraphics[width=#1]{#2}
  \end{center}
  \vspace{0.6cm}
}

% Objekt-auf-Oberflaeche-Skizze fuer die Zeichenaufgabe (LZ2, zweiter Teil):
% Klotz auf einer schraffierten Unterlage, #1 = zusaetzlicher TikZ-Code
% (Kraftpfeile), leer fuer die Aufgabenstellung, gefuellt fuer die Loesung.
\newcommand{\ObjektSkizze}[1]{%
\begin{center}
\begin{tikzpicture}[>=Latex,scale=1.6]
  \draw[thick,darkgray] (-2,0) -- (2,0);
  \foreach \x in {-1.8,-1.4,...,1.8} { \draw[darkgray] (\x,0) -- (\x-0.25,-0.3); }
  \draw[fill=lightgray!60] (-0.5,0) rectangle (0.5,0.7);
  \node at (-0.22,0.35) {$m$};
  #1
\end{tikzpicture}
\end{center}
}

% --- Loesungen ein-/ausblenden -----------------------------------------------
\noprintanswers
%\printanswers

\pointformat{}
\renewcommand{\solutiontitle}{\noindent\color{titleblue}\textbf{L\"osung:}\enspace}

\pagestyle{headandfoot}
\cfoot{\textcolor{darkgray}{Seite \thepage\ von \pageref{LastPage}}}

\newcommand{\Kopfzeile}[4]{%
  \begin{titlebox}
    {\color{white}\sffamily\bfseries\Large #4}
  \end{titlebox}
  \vspace{0.5em}
  \noindent\textbf{Klasse:}~#1 \hfill \textbf{Datum:}~#2 \hfill \textbf{Thema:}~#3
    \hfill \textbf{Lehrperson:}~Loris Keller
  \vspace{0.8em}
  \lhead{\textcolor{darkgray}{#3}}
  \rhead{\textcolor{darkgray}{#4}}
}

\newlength{\antwortraumlen}
\newcommand{\Antwortraum}[1]{%
  \pgfmathsetlength{\antwortraumlen}{#1*2.5cm}%
  \begin{answerbox}
    \rule{0pt}{\antwortraumlen}%
  \end{answerbox}%
}

% Werte fuer Kopfzeile zentral setzen
\newcommand{\Klasse}{2e}
\newcommand{\Datum}{16.09.2026}
\newcommand{\Thema}{Wechselwirkung}
\newcommand{\ArbeitsblattTitel}{Aufgabenblatt 4: Wechselwirkung}

\begin{document}

\Kopfzeile{\Klasse}{\Datum}{\Thema}{\ArbeitsblattTitel}

% =============================================================================
% Aufgabe 1: Zeichenaufgabe 2 -- Objekt auf einer Oberflaeche (LZ2, 2. Teil)
% =============================================================================
\begin{questions}
\begin{taskbox}[Aufgabe 1: Ein Objekt auf einer Oberfläche]
\question[7] Ein Objekt der Masse $m$ liegt auf einer Oberfläche:

\ObjektSkizze{}

\begin{parts}
\part[2] Zeichnen Sie in die Skizze die Kräfte ein, die sich am Objekt
im Gleichgewicht befinden.

\begin{solution}
\ObjektSkizze{
  \draw[->,ultra thick,titleblue] (0,0.35) -- ++(0,1.1) node[midway,right,xshift=2pt]{$\vec F_N$};
  \draw[->,ultra thick,accentorange] (0,0.35) -- ++(0,-1.1) node[midway,right,xshift=2pt]{$\vec F_G$};
}
Gewichtskraft $\vec F_G$ und Normalkraft $\vec F_N$ halten sich am
Objekt exakt die Waage.
\end{solution}

\part[2] Zeichnen Sie stattdessen in eine zweite, identische Skizze das
Wechselwirkungspaar Objekt$\leftrightarrow$Oberfläche ein (an zwei
verschiedenen Körpern: Objekt und Oberfläche).

\ObjektSkizze{}

\begin{solution}
\ObjektSkizze{
  \draw[->,ultra thick,accentpurple] (0,0) -- ++(0,0.9) node[midway,right,xshift=2pt]{$\vec F_N$ (auf Objekt)};
  \draw[->,ultra thick,accentpurple] (0,0) -- ++(0,-0.9) node[midway,right,xshift=2pt]{$-\vec F_N$ (auf Oberfläche)};
}
An der Kontaktstelle drückt die Oberfläche mit $\vec F_N$ nach oben auf
das Objekt, das Objekt drückt nach dem Wechselwirkungsgesetz mit
$-\vec F_N$ gleich stark nach unten auf die Oberfläche.
\end{solution}

\part[3] Das zweite Wechselwirkungspaar dieser Situation ist
Objekt$\leftrightarrow$Erde (zur Gewichtskraft). Beschreiben Sie in
Worten, welche Kraft die Erde dabei erfährt, und wieso sich diese Kraft
in der Praxis nicht bemerkbar macht.

\Antwortraum{2}

\begin{solution}
Nach dem Wechselwirkungsgesetz zieht nicht nur die Erde das Objekt mit
$\vec F_G$ an, sondern das Objekt zieht auch die Erde mit einer exakt
gleich grossen, entgegengesetzten Kraft $-\vec F_G$ an. Weil die Masse
der Erde extrem viel grösser ist als die Masse des Objekts, ist die
dadurch verursachte Beschleunigung der Erde ($a=F/m$ mit riesigem $m$)
völlig unmerklich, genau wie bei der Wand im Skateboardversuch.
\end{solution}
\end{parts}
\end{taskbox}

% =============================================================================
% Aufgabe 2: Zwei weitere Rueckstossprinzip-Beispiele (LZ6)
% =============================================================================
\newpage
\begin{taskbox}[Aufgabe 2: Rückstossprinzip in zwei weiteren Beispielen]
\question[7] Benennen Sie für jedes der folgenden Beispiele das
Wechselwirkungspaar.

\begin{parts}
\part[2] Startendes Flugzeug:

\Bild{5.5cm}{Bilder/flugzeug.jpg}

\Antwortraum{1.5}

\begin{solution}
Die Triebwerke stossen Luft und Abgase nach hinten aus. Luft und Abgase
üben eine gleich grosse Gegenkraft aus, die das Flugzeug nach vorne
beschleunigt.
\end{solution}

\part[2] Eine Person in einem Boot wirft einen Stein:

\Bild{5.5cm}{Bilder/stein_boot.png}

\Antwortraum{1.5}

\begin{solution}
Die Person stösst den Stein nach vorne (in Wurfrichtung). Der Stein übt
dabei eine gleich grosse Gegenkraft auf die Person (und damit das Boot)
aus, wodurch beide in die entgegengesetzte Richtung zurückgestossen
werden.
\end{solution}

\part[3] Vergleichen Sie das Flugzeug-Beispiel mit dem Beispiel
\glqq Person wirft einen Stein aus dem Boot\grqq: In welchem der beiden
Fälle ist der Massenunterschied zwischen den beiden
Wechselwirkungspartnern (Fahrzeug bzw. Person$+$Boot einerseits, die
ausgestossene Masse andererseits) grösser? Was bedeutet das nach
$a=F/m$ für die Rückstossbewegung des jeweils grösseren Partners?

\Antwortraum{2}

\begin{solution}
Beim Flugzeug ist der Massenunterschied zwischen dem Flugzeug selbst und
den ausgestossenen Abgasen extrem gross, die Rückstossbeschleunigung des
Flugzeugs ist deshalb winzig und für das Auge nicht wahrnehmbar
($a=F/m$ mit sehr grossem $m$). Bei Person und Boot ist der
Massenunterschied zum geworfenen Stein viel kleiner, deshalb bewegen
sich Boot und Person merklich zurück. Die Kraft ist in beiden Fällen ein
Wechselwirkungspaar mit gleich grossen Kräften, nur die Grösse der
jeweiligen Masse entscheidet über die Grösse der sichtbaren
Rückstossbewegung.
\end{solution}
\end{parts}
\end{taskbox}

% =============================================================================
% Aufgabe 3: Log-Rolling -- Reflexionsfrage (LZ6)
% =============================================================================
\newpage
\begin{taskbox}[Aufgabe 3: Log-Rolling]
\question[3] Zwei Personen versuchen, sich gegenseitig von einem im
Wasser schwimmenden Baumstamm zu stossen:

\Bild{7cm}{Bilder/logrollingt.jpg}

Auch hier wirken Wechselwirkungspaare zwischen den beiden Personen (über
den Kontakt zwischen Fuss und Stamm und direktes Gegeneinanderdrücken).
Wieso ist das
einzelne Kräftepaar hier deutlich schwerer zu isolieren als bei den
bisherigen Beispielen dieses Themas?

\Antwortraum{2.5}

\begin{solution}
Anders als etwa beim Skateboardversuch oder der Rakete gibt es hier
nicht nur eine einzelne Kontaktstelle mit einem eindeutigen
Kräftepaar. Beide Personen stehen gleichzeitig in mehrfachem Kontakt (mit
dem Stamm und teils direkt miteinander), balancieren dabei fortlaufend
nach, und der Stamm selbst dreht sich unter ihren Füssen mit. Dadurch
wirken mehrere, sich ständig ändernde Kräftepaare gleichzeitig, statt
eines einzigen, klar abgrenzbaren Wechselwirkungspaares wie in den
anderen Beispielen.
\end{solution}
\end{taskbox}

% =============================================================================
% Optional: PhET-Selbstkontrolle (LZ4)
% =============================================================================
\begin{hinweisbox}[Optional: Selbstkontrolle]
Die PhET-Simulation \glqq Forces and Motion: Basics\grqq{}
(Tauziehen-Screen) eignet sich, um die Erklärung aus dem
Skateboardversuch (Kraft vs. Masse vs. Beschleunigung) selbständig noch
einmal zu überprüfen.
\end{hinweisbox}

% =============================================================================
% Zusatzaufgabe (freiwillig): Federwaage-Experiment mit anderen Werten (LZ3)
% =============================================================================
\begin{taskbox}[Zusatzaufgabe (freiwillig): Federwaage mit anderen Massen]
\question[4] Betrachten Sie erneut den Versuchsaufbau mit der Federwaage
zwischen zwei Umlenkrollen aus dem Arbeitsblatt.

\begin{parts}
\part[2] Stellen Sie sich vor, $m_1$ wäre stattdessen
$\SI{300}{\gram}$ (statt $\SI{200}{\gram}$) und hängt weiterhin allein
auf einer Seite, die andere Seite bleibt am Stativ fixiert. Welche
Anzeige erwarten Sie (mit $g\approx\SI{10}{\newton\per\kilogram}$)?

\Antwortraum{1.5}

\begin{solution}
\[
F_G = m\cdot g = \SI{0.3}{\kilogram}\cdot\SI{10}{\newton\per\kilogram}
    = \SI{3}{\newton}
\]
Die Anzeige beträgt $\SI{3}{\newton}$.
\end{solution}

\part[2] Was würde passieren, wenn Sie stattdessen zwei
\emph{unterschiedlich} schwere Massen (z.\,B. $\SI{150}{\gram}$ und
$\SI{50}{\gram}$) an den beiden Rollen hängen liessen, statt wie im
Arbeitsblatt zwei gleich schwere?

\Antwortraum{2}

\begin{solution}
In diesem Fall befindet sich das System nicht mehr im Kräftegleichgewicht:
Die Gewichtskräfte der beiden Massen sind jetzt unterschiedlich gross, die
Zugkraft der (näherungsweise masselosen) Federwaage muss aber weiterhin an
beiden Enden exakt gleich gross sein. Keine der beiden Massen kann dann
noch einzeln für sich im Gleichgewicht sein, das System setzt sich in
Bewegung, statt in Ruhe zu bleiben. Die Anzeige der Federwaage entspricht
dann nicht mehr einfach der Gewichtskraft einer einzelnen Masse, sondern
stellt sich auf einen Wert zwischen den beiden Gewichtskräften ein. Das
ursprüngliche Experiment funktioniert nur deshalb so einfach, weil beide
Massen exakt gleich schwer sind.
\end{solution}
\end{parts}
\end{taskbox}
\end{questions}

\end{document}
